\Article{21}

Sit secundo $n$ quadratum altiorve potestas numeri primi imparis $p$, statuaturque $n=p^{2\varkappa}q$, ita ut sit $q$ vel $=1$ vel $=p$.\quad Hic ante omnia observare convenit, si $\lambda$ sit integer quicunque per $p^\varkappa$ non divisibilis, fieri

\SourceLeaf{43}

\[
\begin{aligned}
&r^{\lambda\lambda}+r^{(\lambda+p^\varkappa q)^2}+r^{(\lambda+2p^\varkappa q)^2}+r^{(\lambda+3p^\varkappa q)^2}\\&\qquad+\text{etc.}+r^{(\lambda+n-p^\varkappa q)^{3}}\\
&\quad=r^{\lambda\lambda}\{1+r^{2\lambda p^\varkappa q}+r^{4\lambda p^\varkappa q}\\&\qquad+r^{6\lambda p^\varkappa q}+\text{etc.}+r^{2\lambda(n-p^\varkappa q)}\}\\&\quad=\frac{r^{\lambda\lambda}(1-r^{2\lambda n})}{1-r^{2\lambda p^\varkappa q}}=0
\end{aligned}
\]

Hinc facile perspicietur, fieri

\[
\begin{aligned}
W=1+r^{p^{2\varkappa}}+r^{4p^{2\varkappa}}+r^{9p^{2\varkappa}}+\text{etc.}+r^{(n-p^\varkappa)^2}
\end{aligned}
\]

Termini enim reliqui progressionis

\[
\begin{aligned}
1+r+r^4+r^9+\text{etc.}+r^{(n-1)^2}
\end{aligned}
\]

distribui poterunt in $(p^\varkappa-1)q$ progressiones partiales, quae singulae sint $p^\varkappa$ terminorum, et per transformationem modo traditam summas evanescentes conficiant.

Hinc colligitur, in casu eo, ubi fit $q=1$, sive ubi $n$ est potestas numeri primi cum exponente pari, fieri

\[
\begin{aligned}
W=p^\varkappa=+\surd n,\;\text{adeoque}\;T=+\surd n,\quad U=0
\end{aligned}
\]

Contra in casu eo, ubi $q=p$, sive ubi $n$ est potestas numeri primi cum exponente impari, statuemus $r^{p^{2\varkappa}}=\rho$, unde $\rho$ erit radix propria aequationis $x^p-1=0$, et quidem $\rho=\cos\frac{k}{p}360^0+i\sin\frac{k}{p}360^0$, ac dein

\[
\begin{aligned}
W&=1+\rho+\rho^4+\rho^9+\text{etc.}+\rho^{(p^{\varkappa+1}-1)^2}\\&=p^\varkappa(1+\rho+\rho^4+\rho^9+\text{etc.}+\rho^{(p-1)^2})
\end{aligned}
\]

Sed summa seriei $1+\rho+\rho^4+\rho^9+\text{etc.}+\rho^{(p-1)^2}$ per art. praec. determinatur, unde sponte concluditur, fieri

\[
\begin{aligned}
W&=\pm\surd n=T,\quad\text{si fuerit }p\text{ formae }4\mu+1\\
W&=\pm i\surd n=iU,\quad\text{si fuerit }p\text{ formae }4\mu+3
\end{aligned}
\]

signo positivo vel negativo valente, prout $k$ fuerit residuum vel non-residuum ipsius $p$.

\Article{22}

Facile quoque ex iis, quae in artt. 20. et 21 exposita sunt, derivatur propositio sequens, quae infra usum notabilem nobis praestabit.\quad Statuatur

\[
\begin{aligned}
W'=1+r^h+r^{4h}+r^{9h}+\text{etc.}+r^{h(n-1)^2}
\end{aligned}
\]

\SourceLeaf{44}

denotante $h$ integrum quemcunque per $p$ non divisibilem, eritque in casu eo, ubi $n=p$, vel ubi $n$ est potestas ipsius $p$ cum exponente impari,

\[
\begin{aligned}
W'&=W,\;\text{si fuerit }h\text{ residuum quadraticum ipsius }p\\
W'&=-W,\;\text{si fuerit }h\text{ non-residuum quadraticum ipsius }p
\end{aligned}
\]

Patet enim, $W'$ oriri ex $W$, si pro $k$ substituatur $kh$; in casu priori autem $k$ et $kh$ similes erunt, in posteriori dissimiles, quatenus sunt residua vel non-residua ipsius $p$.

In casu eo autem, ubi $n$ est potestas ipsius $p$ cum exponente pari, manifesto fit $W'=+\surd n$, adeoque semper $W'=W$.

\Article{23}

In artt. 20. 21. 22 consideravimus numeros primos impares, taliumque potestates: superest itaque casus, ubi $n$ est potestas binarii.

Pro $n=2$ manifesto fit $W=1+r=0$.

Pro $n=4$ prodit $W=1+r+r^4+r^9=2+2r$: hinc $W=2+2i$, quoties $k$ est formae $4\mu+1$, atque $W=2-2i$, quoties $k$ est formae $4\mu+3$.

Pro $n=8$ habemus $W=1+r+r^4+r^9+r^{16}+r^{25}+r^{36}+r^{49}=2+4r+2r^4$ $=4r$.\quad Hinc erit

\[
\begin{aligned}
W&=(1+i)\surd8,\quad\text{quoties }k\text{ est formae }8\mu+1\\
W&=(-1+i)\surd8,\quad\text{quoties }k\text{ est formae }8\mu+3\\
W&=(-1-i)\surd8,\quad\text{quoties }k\text{ est formae }8\mu+5\\
W&=(1-i)\surd8,\quad\text{quoties }k\text{ est formae }8\mu+7
\end{aligned}
\]

Si $n$ est altior potestas binarii, statuamus $n=2^{2\varkappa}q$, ita ut $q$ sit vel $=1$ vel $=2$, atque $\varkappa$ maior quam $1$.\quad Hic ante omnia observari debet, si $\lambda$ sit integer quicunque per $2^{\varkappa-1}$ non divisibilis, fieri

\[
\begin{aligned}
&r^{\lambda\lambda}+r^{(\lambda+2^\varkappa q)^2}+r^{(\lambda+2\pdot2^\varkappa q)^2}+r^{(\lambda+3\pdot2^\varkappa q)^2}\\&\qquad+\text{etc.}+r^{(\lambda+n-2^\varkappa q)^2}\\
&\quad=r^{\lambda\lambda}\{1+r^{2^{\varkappa+1}\lambda q}+r^{2\pdot2^{\varkappa+1}\lambda q}\\&\qquad+r^{3\pdot2^{\varkappa+1}\lambda q}+\text{etc.}+r^{(2n-2^{\varkappa+1}q)\lambda}\}\\&\quad=\frac{r^{\lambda\lambda}(1-r^{2\lambda n})}{1-r^{2^{\varkappa+1}\lambda q}}=0
\end{aligned}
\]

Hinc facile perspicietur, fieri

\[
\begin{aligned}
W=1+r^{2^{2\varkappa-2}}+r^{4\pdot2^{2\varkappa-2}}+r^{9\pdot2^{2\varkappa-2}}+\text{etc.}+r^{(n-2^{\varkappa-1})^2}
\end{aligned}
\]

\SourceLeaf{45}

Statuamus $r^{2^{2\varkappa-2}}=\rho$, eritque $\rho$ radix aequationis $x^{4q}-1=0$, et quidem $\rho=\cos\frac{k}{4q}360^0+i\sin\frac{k}{4q}360^0$; dein fiet

\[
\begin{aligned}
W&=1+\rho+\rho^4+\rho^9+\text{etc.}+\rho^{(2^{\varkappa+1}q-1)^2}\\
&=2^{\varkappa-1}(1+\rho+\rho^4+\rho^9+\text{etc.}+\rho^{(4q-1)^2})
\end{aligned}
\]

Sed summa seriei $1+\rho+\rho^4+\rho^9+\text{etc.}+\rho^{(4q-1)^2}$ per ea, quae de casibus $n=4$, $n=8$ explicavimus, determinatur, unde colligimus

in casu eo, ubi $q=1$, sive ubi $n$ est potestas numeri $4$, fieri

\[
\begin{aligned}
W&=(1+i)2^\varkappa=(1+i)\surd n,\quad\text{si fuerit }k\text{ formae }4\mu+1\\
W&=(1-i)2^\varkappa=(1-i)\surd n,\quad\text{si fuerit }k\text{ formae }4\mu+3
\end{aligned}
\]

quae sunt ipsissimae formulae pro $n=4$ traditae;

in casu eo autem, ubi $q=2$, sive ubi $n$ est potestas binarii cum exponente impari maiori quam $3$, fieri

\[
\begin{aligned}
W&=(1+i)2^\varkappa\surd2=(1+i)\surd n,\quad\text{si fuerit }k\text{ formae }8\mu+1\\
W&=(-1+i)2^\varkappa\surd2=(-1+i)\surd n,\quad\text{si fuerit }k\text{ formae }8\mu+3\\
W&=(-1-i)2^\varkappa\surd2=(-1-i)\surd n,\quad\text{si fuerit }k\text{ formae }8\mu+5\\
W&=(1-i)2^\varkappa\surd2=(1-i)\surd n,\quad\text{si fuerit }k\text{ formae }8\mu+7
\end{aligned}
\]

quae quoque prorsus conveniunt cum iis, quae pro $n=8$ tradidimus.

\Article{24}

Etiam hic operae pretium erit, rationem summae progressionis

\[
\begin{aligned}
W'=1+r^h+r^{4h}+r^{9h}+\text{etc.}+r^{h(n-1)^2}
\end{aligned}
\]

ad $W$ determinare, ubi $h$ integrum quemcunque imparem denotat.\quad Quum $W'$ oriatur ex $W$, mutando $k$ in $kh$, valor ipsius $W'$ perinde a forma numeri $kh$ pendebit, ut $W$ a forma ipsius $k$.\quad Statuamus $\frac{W'}{W}=l$, patetque

I.\quad in casu eo, ubi $n=4$, vel altior potestas binarii cum exponente pari, fieri

\[
\begin{aligned}
l&=1,\quad\text{si fuerit }h\text{ formae }4\mu+1\\
l&=-i,\quad\text{si fuerit }h\text{ formae }4\mu+3,\;\text{atque }k\text{ formae }4\mu+1\\
l&=+i,\quad\text{si fuerit }h\text{ formae }4\mu+3,\;\text{atque }k\text{ eiusdem formae.}
\end{aligned}
\]

\SourceLeaf{46}

II.\quad in casu eo, ubi $n=8$, vel altior potestas binarii cum exponente impari, fieri

\[
\begin{aligned}
l&=1,\ \text{si fuerit}\  h\text{ formae }8\mu+1,\\
l&=-1,\ \text{si fuerit}\  h\text{ formae }8\mu+5,\\
l&=+i,\ \text{si fuerit vel }h\text{ formae }8\mu+3,\;\text{atque }k\text{ formae }4\mu+1,\\
&\phantom{=+i,}\text{vel }h\text{ formae }8\mu+7,\;\text{atque }k\text{ formae }4\mu+3,\\
l&=-i,\ \text{si fuerit vel }h\text{ formae }8\mu+3,\;\text{atque }k\text{ formae }4\mu+3,\\
&\phantom{=+i,}\text{vel }h\text{ formae }8\mu+7,\;\text{atque }k\text{ formae }4\mu+1.
\end{aligned}
\]

Per praecc. determinatio summae $W$ pro iis casibus, ubi $n$ est numerus primus vel numeri primi potestas, complete perfecta est: superest itaque, ut eos quoque casus absolvamus, ubi $n$ e pluribus numeris primis compositus est, huc viam nobis sternet theorema sequens:
